\extrachap{Conventions}

In the whole book, we adopt the following conventions:
\begin{itemize}
    \item A complex space is always assumed to be \emph{reduced}, \emph{paracompact} and \emph{Hausdorff}.
    \item A \emph{modification} of a complex space $X$ is a proper bimeromorphic morphism $\pi\colon Y\rightarrow X$ that is locally obtained from a finite composition of blow-ups with smooth centers.
    \item A \emph{subnet} of a net refers to a Willard subnet. More precisely, if $f\colon I\rightarrow X$ and $f'\colon I'\rightarrow X$ are nets, namely, maps from directed sets $I$ and $I'$ to the same set $X$, we say $f'$ is a subnet of $f$ if there is an \emph{increasing} final function $h\colon I'\rightarrow I$ such that $f\circ h=f'$.
    \item A \emph{domain} in $\mathbb{R}^d$ or $\mathbb{C}^n$ refers to a connected open subset.
    \item A \emph{complex manifold} is assumed to be paracompact.
    \item A \emph{submanifold} of a complex manifold means a closed complex submanifold.
    \item A \emph{neighborhood} is not necessarily open.
    \item The set $\mathbb{N}$ of natural numbers includes $0$.
    \item \emph{Increasing functions} and \emph{decreasing functions} are not necessarily strictly monotone.
    \item \emph{Positive forms} are not necessarily strictly positive.
    \item The function $\mathbb{R}\ni x \mapsto x^2$ is \emph{convex} instead of concave. This is different from the convention in \cite{CLS11}.
    \item We say a positive measure $\mu$ on a measurable space $X$ is \emph{carried} by a measurable subset $Y$ if $\mu(X\setminus Y)=0$.
    \item $\ddc$ means $(2\pi)^{-1}\mathrm{i}\partial \overline{\partial}$.
\end{itemize}

We also use the following notation throughout the book.
\begin{itemize}
    \item The indicator function of a set $E$ is denoted by $\mathbf1_E$.
    \item For subsets $A$ and $B$ of a topological space, $A\Subset B$ means that the closure of $A$ is a compact subset of $B$; we then say that $A$ is \emph{relatively compact} in $B$.
    \item The closure, interior and boundary of a set $E$ in its ambient space are denoted by $\overline E$, $\Int E$ and $\partial E$. The support of a function, measure, distribution or current is denoted by $\Supp$.
    \item For an open set $\Omega$ in $\mathbb R^d$ or in a manifold, $C^k(\Omega)$ is the space of $k$ times continuously differentiable functions, where $0\leq k\leq\infty$, and $C^k_c(\Omega)$ is its subspace of functions with compact support. We write $L^p_{\loc}(\Omega)$ for the functions whose restrictions to compact subsets of $\Omega$ are $p$-integrable.
    \item The imaginary unit is $\mathrm{i}$, and $\Real z$ and $\Img z$ are the real and imaginary parts of a complex number $z$.
\end{itemize}
